% ------------------------------------------------------------------
\begin{frame}{今日のゴール / Today's Goal}
  \begin{block}{目標}
    比例フィードバック制御を実装し、初の制御飛行を行う
  \end{block}

  \vspace{1em}

  \begin{itemize}
    \item 閉ループ制御の考え方
    \item P 制御で角速度を安定化
    \item ARM/DISARM による安全管理
  \end{itemize}
\end{frame}

% ------------------------------------------------------------------
\begin{frame}{閉ループ制御 / Closed-Loop Control}
  \begin{center}
    \resizebox{0.85\textwidth}{!}{\documentclass[border=10pt]{standalone}
\usepackage{tikz}
\usetikzlibrary{arrows.meta, positioning, shapes.geometric, calc}

\begin{document}

\definecolor{blockcolor}{RGB}{0, 120, 200}
\definecolor{arrowcolor}{RGB}{80, 80, 80}
\definecolor{fbcolor}{RGB}{220, 50, 30}
\definecolor{sumcolor}{RGB}{40, 160, 40}
\begin{tikzpicture}[
    block/.style={
      rectangle, draw=blockcolor, line width=1.5pt,
      fill=blockcolor!8, rounded corners=3pt,
      minimum width=26mm, minimum height=14mm,
      font=\bfseries, text=blockcolor, align=center
    },
    sum/.style={
      circle, draw=sumcolor, line width=1.5pt,
      fill=sumcolor!8, minimum size=10mm,
      font=\large\bfseries, text=sumcolor
    },
    arrow/.style={-{Stealth[length=3mm]}, line width=1.2pt, arrowcolor},
    lbl/.style={font=\small, above, text=arrowcolor},
  ]

  % Blocks
  \node[block] (target) at (0, 0) {Target\\Rate};
  \node[sum]   (sum)    at (3, 0) {$\Sigma$};
  \node[block] (kp)     at (6, 0) {P Control\\$K_p \times e$};
  \node[block] (mixer)  at (9.5, 0) {Motor\\Mixer};
  \node[block] (drone)  at (13, 0) {StampFly\\Drone};

  % Forward path
  \draw[arrow] (target) -- node[lbl] {$r$} (sum);
  \draw[arrow] (sum) -- node[lbl] {$e$} (kp);
  \draw[arrow] (kp) -- node[lbl] {$u$} (mixer);
  \draw[arrow] (mixer) -- (drone);

  % Output
  \draw[arrow] (drone.east) -- ++(1.5, 0) node[lbl, right] {response};

  % Feedback path
  \node[block, fill=fbcolor!8, draw=fbcolor, text=fbcolor]
    (gyro) at (8, -3) {Gyro\\Sensor};
  \draw[arrow, fbcolor] (13, 0) -- (13, -3) -- (gyro);
  \draw[arrow, fbcolor] (gyro) -- (3, -3) -- (sum);

  % Sum signs
  \node[font=\scriptsize, sumcolor, anchor=south east] at ($(sum.north)+(150:0.4)$) {$+$};
  \node[font=\scriptsize, fbcolor, anchor=north east] at ($(sum.south)+(210:0.4)$) {$-$};

  % Labels
  \node[font=\small, arrowcolor, above] at ($(drone.east)+(0.75, 0.5)$) {angular rate};
  \node[font=\small, fbcolor, below] at (8, -3.8) {feedback};

  % Title
  \node[font=\Large\bfseries, above=15mm of mixer]
    {Closed-Loop Control};

\end{tikzpicture}
\end{document}
}
  \end{center}
\end{frame}

% ------------------------------------------------------------------
\begin{frame}{制御 API}
  \begin{table}
    \centering
    \begin{tabular}{lll}
      \hline
      \textbf{関数} & \textbf{説明} & \textbf{値域} \\
      \hline
      \api{gyro\_x/y/z()} & ジャイロ角速度 & rad/s \\
      \api{rc\_roll()} & ロールスティック & $-1.0 \sim +1.0$ \\
      \api{rc\_pitch()} & ピッチスティック & $-1.0 \sim +1.0$ \\
      \api{rc\_yaw()} & ヨースティック & $-1.0 \sim +1.0$ \\
      \api{rc\_throttle()} & スロットル & $0.0 \sim 1.0$ \\
      \api{motor\_mixer(T,R,P,Y)} & モーターミキサー & --- \\
      \api{is\_armed()} & ARM 状態確認 & true/false \\
      \hline
    \end{tabular}
  \end{table}
\end{frame}

% ------------------------------------------------------------------
\begin{frame}[fragile]{実習: 3軸 P 制御 / Hands-on}
  \begin{lstlisting}[style=sfcpp, title={user\_code.cpp}]
#include "workshop_api.hpp"
void setup() { ws::print("Lesson 5: Rate P-Control"); }
void loop_400Hz(float dt) {
    if (!ws::is_armed()) {
        ws::motor_stop_all();
        ws::led_color(50, 0, 0);  // Red = DISARM
        return;
    }
    ws::led_color(0, 50, 0);  // Green = ARM
    float Kp_rp = 0.5f, Kp_yaw = 2.0f;
    float rate_max = 1.0f, yaw_max = 5.0f;
    float roll_e  = ws::rc_roll()*rate_max  - ws::gyro_x();
    float pitch_e = ws::rc_pitch()*rate_max - ws::gyro_y();
    float yaw_e   = ws::rc_yaw()*yaw_max   - ws::gyro_z();
    ws::motor_mixer(ws::rc_throttle(),
        Kp_rp*roll_e, Kp_rp*pitch_e, Kp_yaw*yaw_e);
}
  \end{lstlisting}
\end{frame}

% ------------------------------------------------------------------
\begin{frame}{安全注意 / Safety}
  \begin{alertblock}{フライト前チェック}
    \begin{itemize}
      \item[$\square$] \warn{保護メガネを着用}
      \item[$\square$] \warn{低スロットルから徐々に上げる}（最初は 30\% 以下）
      \item[$\square$] プロペラの回転方向を確認（M1=CCW, M2=CW, M3=CCW, M4=CW）
      \item[$\square$] 異常時はすぐにスロットルを下げて DISARM
    \end{itemize}
  \end{alertblock}
\end{frame}

% ------------------------------------------------------------------
\begin{frame}{チェックポイント / Checkpoint}
  \begin{alertblock}{確認事項}
    \begin{itemize}
      \item[$\square$] ARM 後スロットルを上げて安定ホバリング
      \item[$\square$] スティック操作で機体が応答する
      \item[$\square$] DISARM で即座にモーターが停止する
    \end{itemize}
  \end{alertblock}

  \vspace{1em}

  \begin{block}{次のレッスン}
    Lesson 6: システムモデリング --- 伝達関数からゲインを設計
  \end{block}
\end{frame}
